% Lab 3 — Generating Reports（原书 PDF p.33–50）
\bichapter{Lab 3：生成报表}{Lab 3: Generating Reports}

\begin{objectiveBox}
完成本实验后，你应能够：
\begin{itemize}
  \item 为 ORCA 中的违例生成汇总报告；
  \item 分析 setup 与 hold 时序报告；
  \item 使用正确的时序报告开关（switches）；
  \item 识别半时钟周期（half clock cycle）路径。
\end{itemize}
\end{objectiveBox}

\begin{durationBox}
约 \textbf{45} 分钟。
\end{durationBox}

\section{概述}
\label{lab3:overview}

\noindent\textit{Overview}

本实验大致步骤如下：
\begin{enumerate}
  \item 恢复 PrimeTime 会话；
  \item 生成各类汇总报告；
  \item 生成并分析时序报告。
\end{enumerate}

\subsection{相关文件与目录}
\label{lab3:files}

\noindent\textit{Relevant Files and Directories}

本实验全部文件位于主目录下的 \cmd{lab3_reports} 目录中
（该目录为指向 \cmd{lab2_constraints} 的符号链接）。

\begin{tabularx}{\textwidth}{@{}l X@{}}
\toprule
\textbf{路径 / 文件} & \textbf{说明} \\
\midrule
\cmd{lab3_reports/} & 当前工作目录 \\
\cmd{orca_savesession/} & 用于恢复实验的会话目录 \\
\cmd{.synopsys_pt.setup} & 自动读取的 PT setup 文件 \\
\bottomrule
\end{tabularx}


\section{操作说明}
\label{lab3:instructions}

\noindent\textit{Instructions}

\subsection{任务 1：为 Lab 3 配置 PrimeTime}
\label{lab3:task1}

\noindent\textit{Task 1. Setup PrimeTime For Lab 3}

\begin{enumerate}
  \item 从 \cmd{lab3_reports} 目录（指向 \cmd{lab2_constraints} 的 workshop 实验目录）启动 PrimeTime，
        并使用 \cmd{orca_savesession} 目录恢复会话：
\begin{lstlisting}
unix% cd lab3_reports
unix% pt_shell
pt_shell> restore_session orca_savesession
\end{lstlisting}

  \item 查找控制许多报告有效数字位数的变量，并将其设为 4 位有效数字。
        \textbf{［提示：aa significant］}
\begin{lstlisting}
pt_shell> set report_default_significant_digits 4
\end{lstlisting}

\begin{questionBox}
\textbf{问题 1.}\ 该变量的默认值是多少？\textbf{［提示：man page］}\\[0.35em]
\textbf{答：}\\
\begin{lstlisting}
pt_shell> set_app_var report_default_significant_digits -default
> report_default_significant_digits = "2"

pt_shell> set report_default_significant_digits 4
\end{lstlisting}
默认值为 \textbf{2} 位有效数字。
\end{questionBox}
\end{enumerate}

\subsection{任务 2：生成汇总报告}
\label{lab3:task2}

\noindent\textit{Task 2. Generate Summary Reports}

从 Lab~1 可知 ORCA 中存在 setup 违例。

\begin{enumerate}
  \item 通过生成合适的汇总报告回答下列问题：

\begin{questionBox}
\textbf{问题 2.}\ 找出 slack 最差的 5 个 setup 违例，需给出端点名称与 slack。\\[0.35em]
\textbf{答：} 下列命令将列出按 slack 排序的全部 setup 违例。
使用 page mode 可从长报告中退出——所需信息仅为前 5 个违例。
\begin{lstlisting}
pt_shell> report_analysis -status violated -check setup -nosplit
Constrained Pin                          Related Pin  Clock    Type   Slack
I_ORCA_TOP/I_BLENDER/s4_op2_reg[31]/D      CP(rise)     SYS_CLK  setup  -0.9872
I_ORCA_TOP/I_BLENDER/s4_op1_reg[31]/D      CP(rise)     SYS_CLK  setup  -0.8410
I_ORCA_TOP/I_BLENDER/s4_op2_reg[30]/D      CP(rise)     SYS_CLK  setup  -0.8305
I_ORCA_TOP/I_BLENDER/s4_op1_reg[15]/D      CP(rise)     SYS_CLK  setup  -0.6918
I_ORCA_TOP/I_BLENDER/s4_op1_reg[30]/D      CP(rise)     SYS_CLK  setup  -0.6843
\end{lstlisting}

\textbf{问题 3.}\ 列出存在违例 setup 时序路径的 2 个时钟域，以及存在违例 hold 时序路径的 5 个时钟域
（ORCA 共有 6 个时钟域）。\\[0.35em]
\textbf{答：}\\
\begin{lstlisting}
pt_shell> report_constraint -all_violators -max_delay -min_delay
max_delay/setup ('PCI_CLK' group)
max_delay/setup ('SYS_CLK' group)
min_delay/hold ('PCI_CLK' group)
min_delay/hold ('SDRAM_CLK' group)
min_delay/hold ('SD_DDR_CLK' group)
min_delay/hold ('SYS_2x_CLK' group)
min_delay/hold ('SYS_CLK' group)
\end{lstlisting}
违例 setup 的 2 个时钟域：\textbf{PCI\_CLK}、\textbf{SYS\_CLK}。
违例 hold 的 5 个时钟域：\textbf{PCI\_CLK}、\textbf{SDRAM\_CLK}、\textbf{SD\_DDR\_CLK}、\textbf{SYS\_2x\_CLK}、\textbf{SYS\_CLK}。

\textbf{问题 4.}\ 分别有多少 hold 违例位于 input 路径、output 路径，以及 register-to-register 路径？\\[0.35em]
\textbf{答：}\\
\begin{lstlisting}
pt_shell> report_global_timing
Hold violations
                Total   reg->reg  in->reg  reg->out  in->out
WNS            -0.4375  -0.1420   -0.2363  -0.1281   -0.4375
TNS           -12.6317  -3.5616   -2.5768  -0.5954   -5.8979
NUM               100      59       18        7       16
\end{lstlisting}
input 路径 \textbf{18} 个，output 路径 \textbf{7} 个，register-to-register 路径 \textbf{59} 个。
\end{questionBox}

  \item 为以 output 端口 \cmd{sd_DQ[0]} 至 \cmd{sd_DQ[15]} 为终点的 16 位总线，
        生成各 bit 的 setup 最差 slack 报告（这些 output 端口均由单一时钟 \cmd{SD_DDR_CLK} 约束）。

\begin{questionBox}
\textbf{问题 5.}\ 列出 slack 裕量最大（最好 slack）的端点。\\[0.35em]
\textbf{答：} output 端口 \cmd{sd_DQ[0]} 的 slack 最好，为 \textbf{1.5994\,ns}。

一般情况下，下列命令对每个端点仅生成一条报告（\opt{nworst} 默认为 1，且所有 output 端口由同一时钟约束）。
但增大 \opt{max\_paths} 会隐式使用 \opt{slack\_lesser\_than 0}，
而所有到达该端点的 slack 均为正值，因此须将 \opt{slack\_lesser\_than} 改为较大正数（此处为 100）：
\begin{lstlisting}
pt_shell> report_timing -path_type end -max_paths 16 \
    -slack_lesser_than 100 -to sd_DQ*
Endpoint          Path Delay  Path Required  CRP     Slack
sd_DQ[14] (inout)  8.4133 f    9.6649       0.2418  1.4933
...
sd_DQ[0] (inout)   8.0654 f    9.6649       0.0000  1.5994
\end{lstlisting}
\end{questionBox}

  \item 生成设计质量的高层概览：

\begin{questionBox}
\textbf{问题 6.}\ 哪个 clock group 的违例路径数量最多？\\[0.35em]
\textbf{答：}\\
\begin{lstlisting}
pt_shell> report_qor -only violated
\end{lstlisting}
\textbf{SDRAM\_CLK} 有 \textbf{30} 条违例路径，为最多。
\end{questionBox}
\end{enumerate}

\subsection{任务 3：分析 Setup 与 Hold 时序报告}
\label{lab3:task3}

\noindent\textit{Task 3. Analyze Timing Reports for Setup and Hold}

\begin{enumerate}
  \item 打开 page mode：
\begin{lstlisting}
pt_shell> page_on
\end{lstlisting}

  \item 执行下列命令，为 \cmd{PCI_CLK} 生成时序报告：
\begin{lstlisting}
pt_shell> report_timing -group PCI_CLK
\end{lstlisting}

\begin{questionBox}
\textbf{问题 7.}\ 该时序路径是满足（meet）还是违例（violate）时序？\\[0.35em]
\textbf{答：} 违例时序，slack 为 \textbf{$-1.157$\,ns}。

\textbf{问题 8.}\ 该时序路径属于哪一类——内部 flip-flop 到 flip-flop、input，还是 output 时序路径？\\[0.35em]
\textbf{答：} 这是一条 \textbf{output} 时序路径，终点为 output 端口 \cmd{pad[1]}。
\end{questionBox}

  \item 为同一 clock group \cmd{PCI_CLK} 生成 hold 时序报告：
\begin{lstlisting}
pt_shell> report_timing -group PCI_CLK -delay min
\end{lstlisting}

\begin{questionBox}
\textbf{问题 9.}\ 该时序路径属于哪一类——内部 flip-flop 到 flip-flop、input，还是 output 时序路径？\\[0.35em]
\textbf{答：} 这是一条 \textbf{内部}时序路径。终点外观类似 flip-flop，
实为 RAM 时序模型中的众多时序弧之一。

\textbf{问题 10.}\ 该时序路径的数据通路上有多少个单元（cell）？\\[0.35em]
\textbf{答：} 数据通路上仅有 \textbf{1} 个单元——起点 flip-flop。
该时序路径由起点 flip-flop 直接连到终点 flip-flop。

\textbf{问题 11.}\ 起点 flip-flop 的 CP 到 Q pin 使用的是上升沿（rise）单元延迟。
        请给出一个可能的原因：为何这比使用该 flip-flop 更快的下降沿（fall）延迟会导致更差的 hold slack？\\[0.35em]
\textbf{答：} 通常 fall 延迟比 rise 延迟更快，会对 hold 产生更差的 slack！
本例的例外是：若终点 flip-flop 数据 pin 的 fall 跳变导致更小的 library hold time 要求——这正是本例的情况。
\end{questionBox}

  下一步将继续探索并确认你对上述问题的回答。

\begin{questionBox}
\textbf{问题 12.}\ 要确认上述问题的答案，你还需要哪些额外信息？\\[0.35em]
\textbf{答：} 再生成一份以终点下降沿跳变计算数据到达时间的时序报告，并比较两份报告。
由此可确认：终点 flip-flop 数据 pin 的下降跳变确实使 library hold time 更小，从而 slack 更好。
（下一步实验将探索这一点。）
\end{questionBox}

  \item 为同一条时序路径再生成一份 hold 报告，但终点使用下降沿跳变（而非上升沿）。
        请复制粘贴 pin 名以避免拼写错误。
        可参考 ``timing reports'' job aid 查找 \cmd{report_timing} 的合适开关。

\begin{questionBox}
\textbf{问题 13.}\ 你使用了该报告中的哪些行来确认报告了正确的路径？\\[0.35em]
\textbf{答：}\\
\begin{lstlisting}
pt_shell> report_timing -delay_type min_fall \
    -to I_ORCA_TOP/I_PCI_READ_FIFO/PCI_RFIFO_RAM/A1[1] \
    -from I_ORCA_TOP/I_PCI_READ_FIFO/count_int_reg[1]/CP
Startpoint: I_ORCA_TOP/I_PCI_READ_FIFO/count_int_reg[1]
            (rising edge-triggered flip-flop clocked by PCI_CLK)
Endpoint:   I_ORCA_TOP/I_PCI_READ_FIFO/PCI_RFIFO_RAM
            (rising edge-triggered flip-flop clocked by PCI_CLK)
Path Group: PCI_CLK
Path Type:  min
\end{lstlisting}

\textbf{问题 14.}\ 猜测是否正确——更快的 fall 延迟导致更快的数据到达时间，
        但更小的 hold 时间要求，从而得到更好的 slack？\\[0.35em]
\textbf{答：} 是！数据到达时间更快（\textbf{1.139\,ns} 对 \textbf{1.150\,ns}），
但 hold 时间要求更小（\textbf{0.211\,ns} 对 \textbf{0.411\,ns}），因此 slack 更好。
Recall：hold time（以及 setup time）是终点 flip-flop 数据 pin 跳变的函数。
\begin{lstlisting}
data arrival time                                  1.139
library hold time                    0.212         1.025
data required time                                 1.025
slack (MET)                                        0.114
\end{lstlisting}
\end{questionBox}
\end{enumerate}

\subsection{任务 4：使用正确的时序报告开关}
\label{lab3:task4}

\noindent\textit{Task 4. Apply the Correct Timing Report Switches}

\begin{enumerate}
  \item 通过在 PrimeTime 中实验与探索回答下列问题。
        可参考 ``timing reports'' job aid 识别合适的命令与开关。

\begin{questionBox}
\textbf{问题 15.}\ 写出为每个 path group 生成一条 setup 时序报告的命令。\\[0.35em]
\textbf{答：}\\
\begin{lstlisting}
pt_shell> report_timing -group [get_path_group *]
\end{lstlisting}

\textbf{问题 16.}\ 写出为每个存在违例的 path group 生成一条 setup 时序报告的命令。\\[0.35em]
\textbf{答：}\\
\begin{lstlisting}
pt_shell> report_timing -group [get_path_group *] -slack_lesser_than 0
# 交互式使用时，可缩写命令名或开关，或使用 Tab 补全：
pt_shell> report_timing -slack_less 0
\end{lstlisting}

\textbf{问题 17.}\ ORCA 中存在违例时序路径的两个 path group 名称是什么（答案来自上一题的结果）？\\[0.35em]
\textbf{答：} 两个 path group 为 \textbf{PCI\_CLK} 与 \textbf{SYS\_CLK}。

\textbf{问题 18.}\ 写出为受 \cmd{PCI_CLK} 约束的任意 output 端口生成 setup 最差 slack 时序报告的命令。\\[0.35em]
\textbf{答：}\\
\begin{lstlisting}
pt_shell> help all_*
pt_shell> all_outputs -help
pt_shell> report_timing -to [all_outputs -clock PCI_CLK]
# 或另一种等价写法：
pt_shell> report_timing -to [all_outputs] -group PCI_CLK
\end{lstlisting}

\textbf{问题 19.}\ ORCA 中有若干 latch；写出识别这些 latch 数据 pin 的命令。\\[0.35em]
\textbf{答：}\\
\begin{lstlisting}
pt_shell> all_registers -level_sensitive -data_pins
\end{lstlisting}

\textbf{问题 20.}\ 写出为 \cmd{latched_clk_en_reg} latch 的 D pin 生成 hold 时序报告的命令。\\[0.35em]
\textbf{答：} 请复制粘贴以避免拼写错误：
\begin{lstlisting}
pt_shell> report_timing -delay min \
    -to I_ORCA_TOP/I_BLENDER*/latched_clk_en_reg/D
\end{lstlisting}
\end{questionBox}
\end{enumerate}

\subsection{任务 5：识别半时钟周期路径}
\label{lab3:task5}

\noindent\textit{Task 5. Identify Half-Clock Cycle Paths}

时钟 \cmd{SDRAM_CLK} 在 ORCA 中约束了许多半时钟周期路径
（即从下降沿触发 flip-flop 到上升沿触发 flip-flop，或反之）。
这些路径须仔细监控（例如 \cmd{SDRAM_CLK} 的占空比尚未明确定义，或用于分析时钟偏斜）。

\begin{enumerate}
  \item 执行下列命令报告 \cmd{SDRAM_CLK} 的时钟周期，并据此回答问题：
\begin{lstlisting}
pt_shell> report_clock SDRAM_CLK
\end{lstlisting}

\begin{questionBox}
\textbf{问题 21.}\ 给定 waveform 列中第一个数为 \cmd{SDRAM_CLK} 的第一个上升沿、第二个数为下降沿——
该时钟定义了多大的占空比？\\[0.35em]
\cmd{SDRAM\_CLK} 上升沿在 \textbf{0\,ns}，下降沿在 \textbf{3.75\,ns}，周期为 \textbf{7.50\,ns}。
占空比为 \textbf{50\%}。

\textbf{问题 22.}\ 对于由 \cmd{SDRAM_CLK} 上升沿约束到 \cmd{SDRAM_CLK} 下降沿的时序路径，
描述 setup 时序报告中将使用的具体时钟边。\\[0.35em]
\textbf{答：} 对于由 \cmd{SDRAM\_CLK} 上升沿约束到下降沿的时序路径，setup 使用的时钟边为 \textbf{0\,ns} 到 \textbf{3.75\,ns}。

\textbf{问题 23.}\ 对于同一条时序路径，描述 hold 时序检查将使用的具体时钟边。\\[0.35em]
\textbf{答：} 对于同一条时序路径，hold 使用的时钟边为 \textbf{7.5\,ns} 到 \textbf{3.75\,ns}。
\begin{lstlisting}
# 半时钟周期路径报告命令示例（反斜杠为续行符）
# -delay_type min_max 将同时生成 setup 与 hold 各一份报告
report_timing -rise_from [get_clocks SDRAM_CLK] \
    -fall_to [get_clocks SDRAM_CLK] -delay_type min_max
report_timing -fall_from [get_clocks SDRAM_CLK] \
    -rise_to [get_clocks SDRAM_CLK] -delay_type min_max
\end{lstlisting}
\end{questionBox}

  \item 通过为 \cmd{SDRAM_CLK} 约束的半时钟周期时序路径生成合适的时序报告，确认下表中的信息：

\begin{center}
\small
\begin{tabular}{@{}llrllr@{}}
\toprule
\multicolumn{3}{c}{\textbf{Setup}} & \multicolumn{3}{c}{\textbf{Hold}} \\
\cmidrule(lr){1-3}\cmidrule(lr){4-6}
\textbf{Launch} & \textbf{Capture} & \textbf{Slack} & \textbf{Launch} & \textbf{Capture} & \textbf{Slack} \\
\midrule
Fall & Rise 3.75\,ns & 0.680\,ns & Fall & Rise 3.75\,ns & 3.558\,ns \\
Fall & Rise 7.50\,ns & 0.635\,ns & Fall & Rise 3.75\,ns & 3.514\,ns \\
\bottomrule
\end{tabular}
\end{center}

\begin{questionBox}
\textbf{问题 24.}\ 哪个开关可用于为每类半时钟周期时序路径生成最差的 10 条时序报告？\\[0.35em]
\textbf{答：} 在上述命令中添加开关 \opt{-max\_paths 10}。

\textbf{问题 25.}\ 为何 PrimeTime 报告 ``no constrained paths''？
（提示——PrimeTime 使用的选项显示在 \cmd{report_timing} 命令之后。）\\[0.35em]
\opt{-max\_paths} 会隐式设置另一选项 \opt{slack\_lesser\_than 0}；
由于 slack 为正值，因此没有可报告的路径。

\textbf{问题 26.}\ 要报告最差的 10 条时序路径，还必须使用哪个额外选项？\\[0.35em]
\textbf{答：} 由于不存在违例路径可报告，须添加选项：\opt{slack\_lesser\_than 100}。
\end{questionBox}

  \item 退出 PrimeTime：
\begin{lstlisting}
pt_shell> quit
\end{lstlisting}
\end{enumerate}

\noindent 至此 Lab~3 完成。Day-1 课程结束。
